%%%%%%%%%%%%%%%%%%%%%%%%%%%%%%%%%%%%%%%%%%
\section{Materials and Methods}

\subsection{Problem Formulation, Energy Dissipation and Benchmarks}\label{sec:problem}

Let $\Omega=[-1,1]$ and $t\in[0,T]$. We consider
\begin{equation}\label{eq:AC}
u_t = \varepsilon^2 u_{xx} - f(u), \qquad u_x(\pm1,t)=0, \qquad u(x,0)=u_0(x),
\end{equation}
with $f(u)=u^3-u$. The Ginzburg--Landau free energy is $E[u]=\int_{-1}^{1}\bigl(\frac{\varepsilon^2}{2}u_x^2+F(u)\bigr)\,dx$.

\rev{\emph{Initial condition.} The initial condition should satisfy the compatibility condition $u_0'(\pm1)=0$; otherwise $u_x$ is discontinuous at the corners $(x,t)=(\pm1,0)$ and the solution has an initial boundary layer. The condition $u_0(x)=x^2\cos(\pi x)$ that is widely used for the Allen--Cahn PINN benchmark with \emph{periodic} boundary conditions~\cite{Raissi,ref_Kutuk2025} does not satisfy it for the Neumann problem ($u_0'(\mp1)=\pm2$), and it was used with Neumann conditions in the first version of this study. All one-dimensional experiments of the main text use
\begin{equation}\label{eq:IC}
u_0(x)=\tfrac12\bigl(\cos(\pi x)-\cos^2(\pi x)\bigr)=\sin^2\!\bigl(\tfrac{\pi x}{2}\bigr)\cos(\pi x),
\end{equation}
which has the same values at $x=0,\pm\tfrac12,\pm1$ as $x^2\cos(\pi x)$ (namely $0,0,-1$) and $u_0'(\pm1)=0$. Appendix~\ref{app:orig} repeats the B1 experiments of the as-submitted Spectral-PINN and of the classical solver with the original condition and finds the same behavior; the revised network was not run with it, because its exact Neumann basis cannot represent a nonzero wall slope.}

\rev{\emph{Benchmarks.} B1 is $\varepsilon=0.01$, $T=1$ (the original choice). Its solution is dominated by the reaction term: no interface forms before $t\approx3$, and neglecting diffusion altogether (the closed-form solution of $u'=u-u^3$ at every $x$) already gives an $L^2$ error of \N{ode.compat.0.01.l2} at $t=1$ (Table~\ref{tab:forward}), so B1 tests the PDE only weakly. B2 is $\varepsilon=0.1$, $T=1$: the solution is smooth but diffusion matters (reaction-only error \N{ode.compat.0.1.l2}), so a method that solves the PDE must beat that number clearly. B3 is $\varepsilon=0.05$, $T=5$: sharp interfaces ($\max|u_x|\approx14$) form around $t=3$--$5$; a polynomial of degree $K=32$ cannot represent the solution at $T=5$ to better than $L^2\approx\N{floor.b3.K32}$, while $K=64$ can (\N{floor.b3.K64}).}

\begin{Theorem}[Continuous energy dissipation]\label{thm:energy}
Any smooth solution of Equation~\eqref{eq:AC} satisfies $\dfrac{dE}{dt}=-\displaystyle\int_{-1}^{1}u_t^2\,dx\le0$.
\end{Theorem}
\begin{proof}
Differentiating $E$ under the integral sign and integrating $\int u_x u_{xt}\,dx$ by parts (the boundary term vanishes by the Neumann condition) gives $dE/dt=\int(-\varepsilon^2u_{xx}+f(u))u_t\,dx=-\int u_t^2\,dx$, using Equation~\eqref{eq:AC}.
\end{proof}

This theorem is the structural target that the energy-constrained loss (Section~\ref{sec:energy-loss}) attempts to enforce approximately. \rev{Section~\ref{sec:balance} gives its counterpart for a network that does \emph{not} solve the PDE exactly.}

\subsection{Reference Solution and Baseline Methods}\label{sec:cheb-baseline}

\rev{\emph{Reference solution.} All errors are measured against a spectral reference: Chebyshev collocation on the Gauss--Lobatto nodes with the two Neumann rows used to eliminate the wall unknowns, and the stiff integrator Radau (rtol $10^{-11}$) in time. For $T\le1$ it was certified by $N$-refinement, tolerance tightening, a second integrator and an independent ghost-node finite-difference solution, and for B3 by $N$-refinement, with agreement to $\le10^{-8}$ for all cases used here (Appendix~\ref{app:ref}). Predictions and reference are compared on the same 200 equispaced points of $[-1,1]$, with $\|e\|_{L^2}=(\sum_ie_i^2\,\delta x)^{1/2}$, $\delta x=2/199$, and $\|e\|_{L^\infty}=\max_i|e_i|$. The first version used a first-order finite-difference reference whose wall rows were only first-order accurate (they halve the diffusion at the two wall nodes); its error (up to $4\times10^{-3}$ near the walls at early times) was small compared with the PINN errors but not compared with the classical baseline.}

\emph{Classical baseline.} We discretize Equation~\eqref{eq:AC} on the Chebyshev--Gauss--Lobatto nodes $x_j=\cos(j\pi/N)$ using the differentiation matrix $D$ of~\cite{Trefethen}, with a first-order semi-implicit (IMEX) time step (diffusion implicit, reaction $f(u^n)$ explicit) and homogeneous Neumann conditions imposed by boundary bordering; it plays the role of the spectral scheme of~\cite{Shiyapov2026} for the local equation. \rev{We use $N=80$ and $\delta t=5\times10^{-4}$. On B1 and B2 the error is dominated by the time stepping and falls proportionally to $\delta t$ (Table~\ref{tab:forward} gives $\delta t=10^{-4}$ for B1); on B3 it is limited by the resolution $N=80$, so Table~\ref{tab:forward} also gives $N=128$, $\delta t=10^{-4}$. The scheme is an accuracy and cost reference; it does not reproduce the Fast Memory Algorithm of~\cite{Shiyapov2026}.}

\emph{Vanilla PINN.} The second baseline is a space--time PINN, $u(x,t)\approx\mathcal N(x,t;\theta)$, a fully connected network with $\tanh$ activations (three hidden layers of width 64) trained by minimizing $\mathcal L=\mathcal L_{\rm PDE}+\lambda_{\rm IC}\mathcal L_{\rm IC}+\lambda_{\rm BC}\mathcal L_{\rm BC}$ with $\lambda_{\rm IC}=\lambda_{\rm BC}=10$; $\mathcal L_{\rm PDE}$ is the mean-squared residual of Equation~\eqref{eq:AC} on an $80\times80$ grid of collocation points (all derivatives by automatic differentiation), $\mathcal L_{\rm IC}$ fits $u_0$ and $\mathcal L_{\rm BC}$ penalizes $|\mathcal N_x(\pm1,t)|^2$. Training uses Adam (learning rate $10^{-3}$, cosine annealing) for 8000 epochs, the formulation of~\cite{Raissi,Wight,Mattey}. \rev{As a stronger variant of both PINNs, the Adam phase is followed by full-batch L-BFGS (at most 2000 iterations and 2500 loss evaluations, history 50, strong Wolfe line search) on the same loss (``+ L-BFGS''). For B3 the vanilla PINN uses a $60\times120$ collocation grid.}

\subsection{Coefficient-Space PINN (Spectral-PINN)}\label{sec:spectral}

If $u(x)=\sum_{k=0}^{K}c_kT_k(x)$ is a truncated Chebyshev series, its derivative $u'(x)=\sum_{k=0}^{K-1}b_kT_k(x)$ has coefficients given by the backward recurrence
\begin{equation}\label{eq:cheb-recurrence}
b_{K+1}=b_K=0,\qquad b_{k-1}=b_{k+1}+2kc_k\ \ (k=K,\dots,1),\qquad b_0\leftarrow b_0/2.
\end{equation}
This is linear in $\{c_k\}$, so $\mathbf{b}=D_c\mathbf{c}$ with a fixed matrix $D_c$; the second derivative uses $D_c^2$ (verified against analytic derivatives, errors $10^{-11}$--$10^{-13}$ at $K=20$).\label{sec:cheb-coeff}

Following the idea of time-only coefficient networks~\cite{ref_Xia2023}, the network is a function of time alone (two hidden layers of width 64, $\tanh$) that outputs the instantaneous Chebyshev spectrum,
\begin{equation}\label{eq:spectral-pinn}
\mathbf c(t)=\mathcal N(t;\theta)\in\mathbb R^{K+1}, \qquad u(x,t)=\sum_{k=0}^{K}c_k(t)\,T_k(x).
\end{equation}
Spatial derivatives in the residual are exact ($D_c,D_c^2$); the only derivative computed by automatic differentiation is $\dot{\mathbf c}(t)$. The loss is
\begin{equation}\label{eq:loss}
\mathcal L=\mathcal L_{\rm PDE}+\frac{\lambda_{\rm IC}}{K+1}\,|\mathbf c(0)-\mathbf c_0|^2+\lambda_{\rm BC}\mathcal L_{\rm BC}+\lambda_E\mathcal L_E,\qquad
\mathcal L_{\rm PDE}=\Big\langle\tfrac12\|r(\cdot,t)\|_{L^2}^2\Big\rangle_t,
\end{equation}
with $r=u_t-\varepsilon^2u_{xx}+f(u)$, $\lambda_{\rm IC}=10$, $\lambda_{\rm BC}=1$ and $\mathbf c_0$ the coefficient vector of the least-squares fit of $u_0$ by a polynomial of degree $K$ on the spatial grid described below (it differs from the truncated Chebyshev series of $u_0$ by less than $3\times10^{-6}$ for all $K$ used). The time average is over $100\,T$ equispaced collocation times in $[10^{-4},T]$; before the PDE loss is switched on the network is pretrained for 1000 epochs to output $\mathbf c_0$. Training uses Adam (learning rate $10^{-3}$, cosine annealing) for $6000$ epochs (\rev{B1, B2; for B3, 8000 epochs, 250 collocation times and energy checkpoints spaced $0.1$}). \rev{The nonlinearity $f(u)$ is not applied in coefficient space (for $u^3$ this would need a triple product of series of degree $3K$) but} is evaluated on a spatial grid; \rev{two grids are used. \emph{As submitted:} 300 equispaced points, plain mean for the residual and trapezoidal rule for the energy. \emph{Revised:} $N+1$ Gauss--Lobatto points $x_j=\cos(j\pi/N)$, $N=\max(300,4K+8)$, with Clenshaw--Curtis weights, which makes the discrete $\|r\|_{L^2}^2$ and the discrete energy accurate quadratures and keeps the least-squares fit of $\mathbf c_0$ well conditioned (the Vandermonde matrix on 300 equispaced points has condition number $10^{10}$ at $K=128$ and $10^{15}$ for $K\ge192$).}

\rev{\emph{Neumann condition.} Two ways of imposing $u_x(\pm1,t)=0$ are compared. In the \emph{soft} variant of the first version, $\mathcal L_{\rm BC}$ penalizes $u_x(\pm1,t)=\sum_kc_k\,T_k'(\pm1)$, $T_k'(1)=k^2$, $T_k'(-1)=(-1)^{k+1}k^2$, after dividing these vectors by $K^2$; because the loss squares the residual, the effective weight of $u_x(\pm1)^2$ is $\lambda_{\rm BC}K^{-4}$ ($10^{-6}$ at $K=32$, $6\times10^{-8}$ at $K=64$). In the \emph{hard} (revised) variant the network outputs $K-1$ free numbers $\mathbf d=(d_0,\dots,d_{K-2})$, and $\mathbf c=\Phi\mathbf d$ is the coefficient vector of $\sum_{k=0}^{K-2}d_k\varphi_k$ with the classical Neumann basis of the Chebyshev--Galerkin method~\cite{Shen1995,ShenTangWang}
\begin{equation}\label{eq:neumann-basis}
\varphi_k=T_k-\frac{k^2}{(k+2)^2}\,T_{k+2},\qquad k=0,\dots,K-2 .
\end{equation}
Since $T_k'(\pm1)=(\pm1)^{k+1}k^2$, one has $\varphi_k'(\pm1)=0$; the $\varphi_k$ are linearly independent (distinct degrees $k+2$ for $k\ge1$, $\varphi_0=1$) and the polynomials of degree $\le K$ with $p'(\pm1)=0$ form a space of dimension $K-1$ (two independent linear conditions on $\mathbb P_K$), so they span it. The Neumann condition then holds identically, $\mathcal L_{\rm BC}\equiv0$, and the boundary term of the energy balance of Section~\ref{sec:balance} vanishes.}

\subsection{Energy-Constrained Training}\label{sec:energy-loss}

To translate Theorem~\ref{thm:energy} into a training signal we evaluate the discrete energy $E_\theta(t)$ (on the spatial grid above, with the exact spectral $u_x$) at $M+1$ checkpoints $0=\tau_0<\dots<\tau_M=T$, $M=10\,T$ (spacing $0.1$), and add the one-sided penalty
\begin{equation}\label{eq:Lenergy}
\mathcal L_E(\theta)=\sum_{i=0}^{M-1}\big[\max\big(0,\,E_\theta(\tau_{i+1})-E_\theta(\tau_i)\big)\big]^2
\end{equation}
with weight \rev{$\lambda_E$; unless stated otherwise the networks are trained without it ($\lambda_E=0$), and the penalized runs use $\lambda_E=50$, the value of all penalized runs of the first version (Section~\ref{sec:results-energy} varies it from $0$ to $5000$)}. The term is zero whenever the network's energy trajectory is already monotone at the checkpoints; it is a soft analogue of the algebraic energy identity that SAV-type schemes~\cite{Shen_SAV} enforce exactly. \rev{Penalizing energy growth in a PINN loss for Allen--Cahn was proposed in~\cite{ref_Kutuk2025} (with $\partial E/\partial t$ by automatic differentiation at collocation points); our variant penalizes finite increments of the checkpointed energy of a coefficient network. Energy-guided adaptive sampling and training windows are proposed in~\cite{ref_Li2026,Guo_EnergySeq2024}.} \rev{For evaluation, an ``energy violation'' is an increase of the energy by more than $10^{-10}$ between consecutive points of an equispaced time grid. For the Spectral-PINN the energy is computed by Clenshaw--Curtis quadrature on the $4K+9$ Gauss--Lobatto nodes $\cos(j\pi/(4K+8))$, which is exact for the polynomial integrand, and counted for spacings $0.02$, $0.002$ and $2\times10^{-4}$; for the vanilla PINN it is computed by the trapezoidal rule on 1001 equispaced points with $u_x$ by automatic differentiation, and counted for spacing $0.02$ only.}

\subsection{A-Posteriori Energy Balance for a Trained Trajectory}\label{sec:balance}

\rev{Theorem~\ref{thm:energy} holds for exact solutions. A trained network satisfies the PDE only up to the residual $r=u_t-\varepsilon^2u_{xx}+f(u)$ and, in the soft variant, the Neumann condition only up to a defect. The following statement holds for \emph{every} differentiable coefficient trajectory; it uses nothing about the network beyond $\mathbf c\in C^1$. Here $\|\cdot\|$ and $\langle\cdot,\cdot\rangle$ are the norm and inner product of $L^2(-1,1)$.}

\begin{Theorem}[A-posteriori energy balance]\label{thm:balance}
\rev{Let $I\supset[0,T]$ be an open interval, $\mathbf c\in C^1(I;\mathbb R^{K+1})$, $u(x,t)=\sum_{k=0}^Kc_k(t)T_k(x)$, $r=u_t-\varepsilon^2u_{xx}+f(u)$ and $E(t)=\int_{-1}^1\bigl(\frac{\varepsilon^2}{2}u_x^2+F(u)\bigr)dx$. Then $E\in C^1(I)$ and
\begin{align}
\frac{dE}{dt}&=-\|u_t\|^2+\langle r,u_t\rangle+\beta(t) \label{eq:balance-main}\\
&=\tfrac12\|r\|^2-\tfrac12\|u_t\|^2-\tfrac12\|G\|^2+\beta(t)\ \le\ \tfrac12\|r\|^2-\tfrac12\|u_t\|^2+\beta(t),\notag
\end{align}
where $G=-\varepsilon^2u_{xx}+f(u)$ and $\beta=\varepsilon^2\bigl[u_xu_t\bigr]_{x=-1}^{x=1}$ is the boundary defect. Consequently, for $0\le s\le t\le T$,
\begin{equation}\label{eq:balance-int}
E(t)+\tfrac12\int_s^t\|u_t\|^2\,d\tau\ \le\ E(s)+\tfrac12\int_s^t\|r\|^2\,d\tau+\int_s^t\beta\,d\tau .
\end{equation}}
\end{Theorem}
\begin{proof}
\rev{\emph{(1) Differentiability.} For $a\in\mathbb R^{K+1}$ let $u_a=\sum_ka_kT_k$ and $\mathcal E(a)=\int_{-1}^1\bigl(\frac{\varepsilon^2}{2}(u_a')^2+\frac14u_a^4-\frac12u_a^2+\frac14\bigr)dx$. The integrand is a finite sum $\sum_\alpha a^\alpha q_\alpha(x)$ of monomials in $a$ with polynomial coefficients $q_\alpha$, so $\mathcal E$ is a polynomial in $a$, and $E(t)=\mathcal E(\mathbf c(t))$ is $C^1$ with $E'=D\mathcal E(\mathbf c)[\mathbf c']$.
\emph{(2) Differential.} For $v\in\mathbb R^{K+1}$ and $w=u_v$, the function $s\mapsto\mathcal E(a+sv)$ is a polynomial whose coefficient of $s$ is the integral of the coefficient of $s$ in the integrand; hence $D\mathcal E(a)[v]=\int_{-1}^1(\varepsilon^2u_a'w'+f(u_a)w)\,dx$.
\emph{(3) Integration by parts.} For polynomials $p,q$, $\int_{-1}^1pq'\,dx=[pq]_{-1}^1-\int_{-1}^1p'q\,dx$. With $p=u_a'$, $q=w$,
\[
D\mathcal E(a)[v]=\int_{-1}^1\bigl(-\varepsilon^2u_a''+f(u_a)\bigr)w\,dx+\varepsilon^2[u_a'w]_{-1}^{1}.
\]
For $a=\mathbf c(t)$ and $v=\mathbf c'(t)$ one has $w=u_t$, so $E'=\langle G,u_t\rangle+\beta$.
\emph{(4) Algebra.} By the definition of $r$, $u_t=-G+r$, so $\langle G,u_t\rangle=\langle r-u_t,u_t\rangle=-\|u_t\|^2+\langle r,u_t\rangle$, which is the first line of~\eqref{eq:balance-main}. Also $r=u_t+G$ gives $\|r\|^2=\|u_t\|^2+2\langle G,u_t\rangle+\|G\|^2$, i.e.\ $\langle G,u_t\rangle=\frac12\|r\|^2-\frac12\|u_t\|^2-\frac12\|G\|^2$; dropping $-\frac12\|G\|^2\le0$ gives the inequality. Integrating it over $[s,t]$ (fundamental theorem of calculus for the $C^1$ function $E$) gives~\eqref{eq:balance-int}.}
\end{proof}

\rev{For $r=0$ and $\beta=0$ the identity reduces to Theorem~\ref{thm:energy}. The theorem does not say that a trained network dissipates energy, nor that $r$ is small: it says that the energy of the learned solution can grow only through the residual and the boundary defect, both of which are computable a posteriori, because $r^2$ and the energy integrand are polynomials in $x$ (of degree $6K$ and $4K$). The energy and all terms of the identity are integrated exactly by Clenshaw--Curtis quadrature on $4K+9$ nodes, and $\|r\|^2$ by the same rule to a relative accuracy of about \N{quad.max_rel} for the seed-42 network of each configuration of Section~\ref{sec:results-balance} (Appendix~\ref{app:balance}). With the hard Neumann basis $\beta\equiv0$. We evaluate all terms for every Spectral-PINN trained for the forward problem (Section~\ref{sec:results-balance}).}

\subsection{Inverse Problem: Recovering the Interface Parameter}\label{sec:inverse}

For the inverse problem $\varepsilon$ is unknown and synthetic observations $u_{\rm obs}$ of the reference solution (with known $\varepsilon_{\rm true}$), sampled at $n_x=40$ equispaced points and the five times $t\in\{0.2,0.4,0.6,0.8,1.0\}$ and optionally corrupted by Gaussian noise of standard deviation $\sigma$, are given. The loss is $\mathcal L=\lambda_{\rm PDE}\mathcal L_{\rm PDE}(\varepsilon)+\lambda_{\rm IC}\mathcal L_{\rm IC}+\lambda_{\rm BC}\mathcal L_{\rm BC}+\lambda_{\rm data}\mathcal L_{\rm data}$ with $\mathcal L_{\rm IC}=|\mathbf c(0)-\mathbf c_0|^2/(K+1)$ as in Equation~\eqref{eq:loss}, $\mathcal L_{\rm data}=\langle(u_\theta(x_{\rm obs},t_{\rm obs})-u_{\rm obs})^2\rangle$ and $(\lambda_{\rm PDE},\lambda_{\rm IC},\lambda_{\rm BC},\lambda_{\rm data})=(50,10,1,2)$ unless stated otherwise. Three estimators are compared. \emph{Joint Adam} (the procedure of the first version): $\log\varepsilon$ is an additional trainable scalar updated by Adam with a learning rate $10\times$ that of the network, started at $\varepsilon_{\rm init}$. \rev{\emph{Variable projection:} the residual is linear in $\varepsilon^2$, so for the current network the minimizer of $\mathcal L_{\rm PDE}$ over $\varepsilon^2$ is $\varepsilon^2_{\rm LS}=\bigl\langle\langle u_t+f(u),u_{xx}\rangle\bigr\rangle_t\big/\bigl\langle\|u_{xx}\|^2\bigr\rangle_t$ ($\langle\cdot\rangle_t$: average over the collocation times of Equation~\eqref{eq:loss}); this value (restricted to $[10^{-8},4]$ and relaxed with weight $0.1$ from epoch to epoch) replaces the trainable scalar, after a data-fitting phase of 3000 epochs with the PDE loss off and data weight $1000$. \emph{Classical least squares:} nonlinear least squares for $\log\varepsilon$ (trust-region reflective, from $2\varepsilon_{\rm true}$) with the reference solver of Appendix~\ref{app:ref} at a lower resolution ($N=96$, rtol $10^{-10}$) as forward model; the data come from the same solver at the reference resolution, so the noise-free case is close to an ``inverse crime'' and the noisy runs are the meaningful comparison. The neural and the classical estimator receive identical noise realizations. The joint-Adam runs use $4000$ epochs (learning rate $2\times10^{-3}$, cosine annealing) after the 1000-epoch warm-up; the settings varied in the ablation are listed in Section~\ref{sec:results-inverse}.}

\subsection{Two-Dimensional Symmetric Extension}\label{sec:2d}

To test symmetry preservation directly we consider the periodic domain $[-\pi,\pi]^2$ with the reflection-symmetric initial condition $u_0(x,y)=0.4\cos(2x)\cos(2y)$, invariant under $x\mapsto-x$, $y\mapsto-y$ and $x\leftrightarrow y$. Since the isotropic Laplacian and $f$ commute with these reflections, the exact solution remains symmetric. The Fourier-Spectral-PINN represents $u(x,y,t)=\sum_{m,n=0}^{M}c_{mn}(t)\cos(mx)\cos(ny)$; every basis function is even in $x$ and in $y$, so $u(-x,y,t)=u(x,y,t)$ and $u(x,-y,t)=u(x,y,t)$ hold \emph{by construction}, independent of training quality, and $\Delta[\cos(mx)\cos(ny)]=-(m^2+n^2)\cos(mx)\cos(ny)$. As a control, the vanilla PINN receives the periodic (but not reflection-symmetric) encoding $(\cos x,\sin x,\cos y,\sin y,t)$, so only the reflection symmetry is architecture-dependent. We report $S(t)=\langle(u(x,y,t)-u(-x,y,t))^2\rangle+\langle(u(x,y,t)-u(x,-y,t))^2\rangle$ for both methods on the endpoint-inclusive $49\times49$ grid $x_i=-\pi+ih$, $y_j=-\pi+jh$, $h=2\pi/48$, which the reflection maps onto itself (Appendix~\ref{app:sym}). This experiment is unchanged from the first version and involves no Neumann condition.

\subsection{Implementation and Evaluation Protocol}\label{sec:protocol}

\rev{All networks are multilayer perceptrons with $\tanh$ activations in double precision, trained with Adam; the one-dimensional runs used one CPU thread per run. Most runs were executed concurrently on a server with an Intel Xeon Gold 6430 (32 cores, 64 hardware threads); the 42 runs of a first pass (41 production runs and a short smoke test, listed in the repository) were executed on a workstation with an Intel Core i5-9400F, and repeating three of them on the server changed their $L^2$ errors by at most \N{proto.repro}, small compared with the seed-to-seed spread. Because many runs were concurrent, wall-clock times are not a clean benchmark; the times in Table~\ref{tab:forward} come from separate runs executed one at a time on the (shared) server, together with the classical solver on the same machine. Seeds are $\{42,123,7,2024,99\}$ unless stated otherwise; ``mean~$\pm$~SD'' uses the sample standard deviation ($n-1$). Every ratio of errors that is not attributed to individual seeds is a ratio of seed means, with 95\% percentile bootstrap intervals (20\,000 resamples of the seeds, independently for numerator and denominator). The maximum-bound diagnostic is $\max_{x,t}|u|$ on a $1001\times(100T+1)$ grid. Counts of trained networks given in the text include the four separate timing runs with seed~42 (revised Spectral-PINN with $K=32$ and $K=64$, as-submitted Spectral-PINN with $K=32$, vanilla PINN) and two repetitions of the revised network with seed~42 ($K=32$ and $K=64$) made for the grid comparison of Section~\ref{sec:results-k}; three of these six runs reproduce the original run exactly. Source code, the certified reference and all raw results are in the repository cited in the Data Availability Statement.}
